% BHU PYQ -- Manually transcribed & error-corrected
% Paper: MTB-602 : Linear Algebra
% Exam: B.A./B.Sc. (Hons) Semester VI Examination, 2022-23

\documentclass[11pt,a4paper]{article}
\usepackage[a4paper,top=2.5cm,bottom=2.5cm,left=2.8cm,right=2.8cm,headheight=14pt]{geometry}
\usepackage{amsmath,amssymb}
\usepackage[T1]{fontenc}
\usepackage[utf8]{inputenc}
\usepackage{lmodern}
\usepackage{microtype}
\usepackage{fancyhdr}
\usepackage{enumitem}
\usepackage{hyperref}
\usepackage{lastpage}
\usepackage{parskip}

\hypersetup{colorlinks=true,urlcolor=black,linkcolor=black,
  pdftitle={MTB-602: Linear Algebra -- BHU B.A./B.Sc. Sem VI 2022-23}}
\pagestyle{fancy}\fancyhf{}
\renewcommand{\headrulewidth}{0.4pt}\renewcommand{\footrulewidth}{0.4pt}
\fancyhead[L]{\small   \textit{Banaras Hindu University}}
\fancyhead[R]{\small   \textit{MTB-602: Linear Algebra}}
\fancyfoot[C]{\small Page  \thepage\ of \pageref{LastPage}}


\newlist{parts}{enumerate}{2}
\setlist[parts,1]{label=(\alph*),leftmargin=*,itemsep=5pt,topsep=3pt}
\setlist[parts,2]{label=(\roman*),leftmargin=*,itemsep=3pt,topsep=2pt}

\newcommand{\pts}[1]{\hfill   \textit{[#1]}}


\begin{document}

\begin{center}
  {\large\bfseries BANARAS HINDU UNIVERSITY}\\[2pt]
  {
\normalsize B.A./B.Sc. (Hons) --- Semester VI Examination, 2022-23}\\[6pt]
  \rule{0.8\linewidth}{0.5pt}\\[6pt]
  {\large\bfseries Paper No. MTB-602: Linear Algebra}\\[4pt]
  {
\normalsize Time: 3 Hours \hfill Maximum Marks: 70}
\end{center}
\rule{\linewidth}{0.4pt}
\smallskip



\noindent   \textit{Instructions: Answer   \textbf{five} questions including Question~1 which is compulsory. Figures in right-hand margin indicate marks.}
\bigskip




\noindent   \textbf{Question 1.}   \textit{(Compulsory.)}

\begin{parts}
  \item Let $V_1$ and $V_2$ be subspaces of a vector space $V(F)$. Show that $V_1 + V_2$ is the subspace generated by $V_1 \cup V_2$. \pts{3}
  \item Let $T: V_1(F)  o V_2(F)$ be a linear transformation. Show that $\ker T = \{0\}$ if and only if $T$ is one-one. \pts{2}
  \item Show that the vector space $F[x]$ of polynomials over $F$ is not finite dimensional. \pts{2}
  \item Show that a matrix $A$ and its transpose $A^T$ have the same characteristic polynomial. \pts{2}
  \item Find the rank and nullity of the linear transformation $T: \mathbb{R}^3(\mathbb{R})  o \mathbb{R}^3(\mathbb{R})$ defined by $T(x_1, x_2, x_3) = (x_1 + x_2 + x_3, x_2 + x_3, x_3)$. \pts{3}
  \item If $x$ and $y$ are vectors in an inner product space, show that:
  \[
  \|x + y\|^2 + \|x - y\|^2 = 2(\|x\|^2 + \|y\|^2)
  \] \pts{2}
\end{parts}

\medskip\rule{\linewidth}{0.2pt}

\medskip


\noindent   \textbf{Question 2.}

\begin{parts}
  \item Define a vector space over a field. Can every Abelian group $V$ be made into a vector space by the scalar multiplication $av = 0$ for all $a \in F$, $v \in V$? Justify your answer. \pts{3}
  \item If $V_1$ and $V_2$ are subspaces of $V(F)$, show that $V_1 \cup V_2$ is not a subspace in general. Further show that $V_1 \cup V_2$ is a subspace if and only if $V_1 \subseteq V_2$ or $V_2 \subseteq V_1$. \pts{5}
  \item Prove that any two bases of a finitely generated vector space $V(F)$ have the same number of elements. \pts{6}
\end{parts}

\medskip\rule{\linewidth}{0.2pt}

\medskip


\noindent   \textbf{Question 3.}

\begin{parts}
  \item If $W_1$ and $W_2$ are two subspaces of a finite dimensional vector space $V(F)$, prove that:
  \[
  \dim W_1 + \dim W_2 = \dim(W_1 + W_2) + \dim(W_1 \cap W_2)
  \] \pts{9}
  \item
  
\begin{parts}
    \item If $v_1, v_2, v_3 \in V(F)$ such that $v_1 + v_2 + v_3 = 0$, show that $\{v_1, v_2\}$ spans the same subspace as $\{v_2, v_3\}$. \pts{3}
    \item Show by an example that linear dependence depends not only on the vector space but on the field as well. \pts{2}
  \end{parts}
\end{parts}

\medskip\rule{\linewidth}{0.2pt}

\medskip


\noindent   \textbf{Question 4.}

\begin{parts}
  \item State and prove the rank-nullity theorem. \pts{7}
  \item Define the dual space of a vector space $V(F)$. If $V(F)$ is an $n$-dimensional vector space with basis $\{e_1, e_2, \dots, e_n\}$, show that there exists a uniquely determined basis $\{e_1^*, e_2^*, \dots, e_n^*\}$ of its dual $V^*$ such that $e_i^*(e_j) = \delta_{ij}$ for $j = 1, 2, \dots, n$, where $\delta_{ij} = 0$ if $i 
eq j$ and $\delta_{ij} = 1$ if $i = j$. \pts{7}
\end{parts}

\medskip\rule{\linewidth}{0.2pt}

\medskip


\noindent   \textbf{Question 5.}

\begin{parts}
  \item If the matrix of a linear transformation $T$ of $V_3(\mathbb{C})$ with respect to the basis $\{(1,0,0),(0,1,0),(0,0,1)\}$ is
  \[
  A = 
\begin{pmatrix} 0 & 1 & -1 \\ 0 & 1 & -1 \\ -1 & -1 & 0 \end{pmatrix}
  \]
  find the matrix of $T$ with respect to the basis $\{(1,1,-1),(-1,0,1),(1,2,1)\}$. \pts{8}
  \item State the Cayley-Hamilton theorem. Verify it for the matrix:
  \[
  A = 
\begin{pmatrix} 0 & 1 & 2 \\ 1 & 2 & 3 \\ 3 & 3 & 1 \end{pmatrix}
  \] \pts{6}
\end{parts}

\medskip\rule{\linewidth}{0.2pt}

\medskip


\noindent   \textbf{Question 6.}

\begin{parts}
  \item Show that a square matrix $A$ of order $n$ is diagonalizable if and only if $A$ has $n$ linearly independent eigenvectors. \pts{7}
  \item Find the eigenvalues and eigenvectors of the matrix:
  \[
  A = 
\begin{pmatrix} 1 & 2 & 3 \\ 0 & -4 & 2 \\ 0 & 0 & 7 \end{pmatrix}
  \] \pts{7}
\end{parts}

\medskip\rule{\linewidth}{0.2pt}

\medskip


\noindent   \textbf{Question 7.}

\begin{parts}
  \item If $\{u_1, u_2, \dots, u_n\}$ is any finite orthogonal set in an inner product space $V(F)$ and $u$ is any vector in $V$, show that:
  \[
  \sum_{k=1}^n |(u, u_k)|^2 \le \|u\|^2
  \]
  Furthermore, show that equality holds if and only if $u$ is in the subspace generated by $\{u_1, u_2, \dots, u_n\}$. \pts{8}
  \item Apply the Gram-Schmidt process to vectors $y_1 = (1,0,1)$, $y_2 = (1,0,-1)$, $y_3 = (0,3,4)$ to obtain an orthonormal basis for $\mathbb{R}^3(\mathbb{R})$ with standard inner product. \pts{6}
\end{parts}

\vfill
\rule{\linewidth}{0.4pt}

\begin{center}\small
  \href{https://bhu-exam.vercel.app/}{Click here to visit our website}\\[4pt]
  \href{https://me-aryan.vercel.app/}{Click here to visit our contributor}\\[4pt]
  \href{https://quashan-me.vercel.app/}{Click here to visit our contributor}
\end{center}

\end{document}
